\subsectionlink{量子コンピューターにおけるデコヒーレンスと制御可能性}
すでにいくつかの箇所で示唆したように、量子コンピューターの能力は、計算基底状態(7.5)の非局所的でコヒーレントな重ね合わせの存在に由来する。これらの計算基底状態は、個々の量子ビットの基底状態$|0\rangle$と$|1\rangle$の積状態である（量子コンピューターを制御する何らかの装置によって操作し、測定できるよう、両状態は十分に区別可能でなければならない）。したがって量子計算は、中間的または巨視的に異なる状態の重ね合わせに基づいており、それらは通常、デコヒーレンスにきわめて敏感である。デコヒーレンスは量子コンピューターの最大の敵であり、量子コンピューターの実現に関する研究の、すべてとはいわなくても多くは、デコヒーレンスをいかに制御し、最小化し、「取り消す」かという問題をめぐって進められてきたといってよい。

量子コンピューターのデコヒーレンス問題に対処するため、量子ビットをあらゆる環境との相互作用からできるだけ遮蔽すれば十分だと提案したくなるかもしれない。確かに、量子ビットを周囲から孤立させれば、デコヒーレンスは最小になる。しかし、この方策は量子コンピューターの機能も損なう。すでに見たように、量子計算は、量子ビットにユニタリー操作の列を適用することで実行される。物理的には、そのようなゲート操作は、量子ビット系と外部装置との何らかの相互作用に対応する。その相互作用の具体的な物理的形態は、量子コンピューターの個別の実現方法によって異なる（表7.1参照）。
% 原文 p.302
\hypertarget{srcpage:302}{}\label{srcpage:302}
量子ゲートを確実に実行するためには、量子ビットと制御装置の相互作用が十分に強くなければならない。

したがって、相反する二つの要求に直面する。一方では、望ましくない環境との相互作用を最小にし、量子ビットの重ね合わせ状態を保護するために、量子ビットを孤立させたい。他方では、巨視的装置との制御された相互作用により、量子ビットを操作し（状態準備とゲート操作の形で）、読み出せるように、系を十分に開いておく必要がある。すると、こうした量子ビットと装置の相互作用は、一般に装置とのもつれによって量子ビットにデコヒーレンスを生じさせ、装置の古典的雑音（揺らぎ）によってさらに擾乱が加わる。

したがって、量子コンピューターの設計という困難な課題は、両方の要求を均衡させて満たすことにある。環境によるデコヒーレンス率が非常に低くても、まともに操作できない量子コンピューターは、よく制御できても、環境との相互作用が強すぎて装置が直ちに古典コンピューターへと劣化してしまう量子コンピューターと同じく、役に立たないことは明らかである。要するに、量子ビットは、環境のある構成要素（すなわち装置）とは強く相互作用し、残りの部分（量子ビットを取り囲む、制御されない周囲）とは相互作用しないことが求められる。

たとえば光量子計算は、光子を物理的量子ビットとして使い、偏光の二つの異なる向きで計算基底状態を表す。光子は電荷をもたず、通常、互いに、また物質とも弱くしか相互作用しないため、デコヒーレンスに比較的鈍感である。通常の線形光学媒質では、光子は互いにまったく相互作用せず（非線形媒質でも、相互作用はきわめて弱い）、そのため、一対の光子をもつれさせる二量子ビットゲート操作をどう実行するかという問題が生じる。このジレンマから抜け出す可能な方法として、有名なKnill--Laflamme--Milburn方式\cite{ref348}が提案されている。この方式は、線形光学、単一光子状態、射影測定だけを用いて、二つの光子量子ビットをもつれさせる。しかし今度は、この方式自体が、光量子コンピューターの構築を可能にするために克服すべき新たな課題を生む。同様に、原子核内の単一スピンとして実現した量子ビットは、デコヒーレンス時間が長いが、制御と測定が非常に難しい。

\subsectionlink{デコヒーレンスと古典的揺らぎ}
2.12節で論じたように、デコヒーレンスは、古典的な対応物のない、明確に量子力学的な効果として理解すべきである。これに対して量子計算の文献では、「デコヒーレンス」という用語が、古典的な揺らぎや不完全さによる擾乱を含め、量子ビットに影響するあらゆる過程を指す、もっと大まかな意味で使われていることが多い。
% 原文 p.303
\hypertarget{srcpage:303}{}\label{srcpage:303}
量子計算におけるこのような古典的雑音の源には、イオントラップ内の量子ビットに入射するレーザー光の強度\cite{ref349}と持続時間\cite{ref350}の揺らぎ、NMR量子計算での磁場の不均一性\cite{ref351}、超伝導量子ビットでのバイアス揺らぎ\cite{ref293}などがある。

現象論的にも形式的にも、このような古典的過程が量子ビットに及ぼす影響は、環境とのもつれの効果と同様に、すなわち「量子ビット誤り」によって記述できる（後の7.4.2節参照）。したがって純粋に実用的な観点では、一方にある環境とのもつれを生む相互作用によるデコヒーレンスと、他方にある古典的雑音による（アンサンブル平均での）位相コヒーレンスの喪失との区別は、関係しないかもしれず、しばしば明示されない。量子誤り訂正（7.4節参照）の分野も、誤り訂正方式が量子ビット誤りの物理的起源に依存しないため、古典的雑音と量子デコヒーレンスの区別をさらに曖昧にしている。

さらに2.12節で述べたように、環境とのもつれによる位相コヒーレンスの喪失は、ときに、系を乱す古典的揺らぎ、すなわち系の自己Hamiltonianにある時間依存項を加えることで模擬される。この方法は、たとえば、イオントラップ量子コンピューターの量子ビットの位相コヒーレンスに、揺らぐパラメーターが及ぼす影響の理論的研究\cite{ref349,ref352}と実験的研究\cite{ref125,ref353}で実行された。理論面では、SchneiderとMilburnが、量子ビットを制御するレーザー光の位相と強度のランダムな揺らぎ\cite{ref349}、およびイオンを閉じ込めるトラップのパラメーターの揺らぎによる「デコヒーレンス」を調べた。NISTのグループによる実験\cite{ref125,ref353}では、トラップされたイオンのSchrödingerの猫状態（コヒーレント状態および数固有状態の重ね合わせ）の減衰に対する、異なる種類の熱浴の影響を、調和トラップの極小位置やトラップ周波数などのパラメーターをランダムに揺らがせることで、実験的に模擬した。

2.12節ですでに論じたように、ここで鍵となるのは、環境の影響をこのように古典的雑音として模型化すると、コヒーレンスの減衰はアンサンブルにしか、すなわち、その雑音過程の多数の個別の実現についての平均としてしか現れないことである。系は外部環境に結合していないので、雑音過程の個々の実現では系の力学は完全にユニタリーであり、系からコヒーレンスが失われているはずはない。これに対して、系が環境の自由度ともつれる場合、適切な逆変換を加えてデコヒーレンスの効果を反転させるのに十分な情報を得るには、少なくとも系との相互作用の前後に、環境に対して一対の測定を行う必要がある。さらにNISTの実験\cite{ref125}で示されたように、これらの測定は、デコヒーレンスを「取り消す」のに、つねに十分な手続きになるわけではない（文献\cite{ref016}のIV.C節も参照）。
% 原文 p.304
\hypertarget{srcpage:304}{}\label{srcpage:304}
したがって、この実験は、古典的雑音と環境によるデコヒーレンスの主要な類似点と相違点のいくつかを、見事に示している。

もちろん、これは古典的雑音が量子計算で重要でないという意味ではない。逆に、実験的実現での揺らぎと不完全さに起因する雑音は、環境とのもつれが生じさせるデコヒーレンスの効果と同じくらい重要かもしれない。以下でも「デコヒーレンス」は環境とのもつれの帰結を記述する用語として用いるが、より一般的な「量子ビット誤り」は、特に断らないかぎり、古典的雑音によるか環境とのもつれによるかを問わず、量子ビット状態のあらゆる変化を指すものとする。
